\documentclass[UTF8,a4paper,zihao=-4,fontset=none]{ctexart}
\newcommand{\ReportHeadLeft}{端到端赛道 ｜ AI 工具使用说明}
\newcommand{\ReportHeadRight}{2026 年 11 月}
% 决赛策略研究报告 公共导言（XeLaTeX + ctexart）
% 原则：层级清楚（章、节两级为主，最多到小节）；正文成段；不用彩色框、侧栏、底纹、小字注释。
% 图表标题在上，写成「图 1　标题（数据范围）」，数据来源写进标题括号或正文，不另起小字行。
%
% 用法：
%   \documentclass[UTF8,a4paper,zihao=-4,fontset=none]{ctexart}
%   % 决赛策略研究报告 公共导言（XeLaTeX + ctexart）
% 原则：层级清楚（章、节两级为主，最多到小节）；正文成段；不用彩色框、侧栏、底纹、小字注释。
% 图表标题在上，写成「图 1　标题（数据范围）」，数据来源写进标题括号或正文，不另起小字行。
%
% 用法：
%   \documentclass[UTF8,a4paper,zihao=-4,fontset=none]{ctexart}
%   % 决赛策略研究报告 公共导言（XeLaTeX + ctexart）
% 原则：层级清楚（章、节两级为主，最多到小节）；正文成段；不用彩色框、侧栏、底纹、小字注释。
% 图表标题在上，写成「图 1　标题（数据范围）」，数据来源写进标题括号或正文，不另起小字行。
%
% 用法：
%   \documentclass[UTF8,a4paper,zihao=-4,fontset=none]{ctexart}
%   \input{../_style/report_preamble.tex}
%   \title{…} \subtitle{…} \team{…} \reportdate{…}
%   \begin{document} \makereporttitle ...

\setCJKmainfont{Songti SC}[BoldFont={Songti SC Bold}]
\def\yhpath{/Applications/Microsoft PowerPoint.app/Contents/Resources/DFonts/}
\setCJKsansfont{msyh.ttc}[Path=\yhpath,BoldFont=msyhbd.ttc]
\setCJKmonofont{msyh.ttc}[Path=\yhpath]
\setmainfont{Times New Roman}
\setsansfont{msyh.ttc}[Path=\yhpath,BoldFont=msyhbd.ttc]
\setmonofont{Menlo}[Scale=0.88]

\usepackage{newunicodechar}
\newunicodechar{→}{\ensuremath{\rightarrow}}
\newunicodechar{≈}{\ensuremath{\approx}}
\newunicodechar{≤}{\ensuremath{\leq}}
\newunicodechar{≥}{\ensuremath{\geq}}
\newunicodechar{−}{\ensuremath{-}}
\newunicodechar{×}{\ensuremath{\times}}
\newunicodechar{÷}{\ensuremath{\div}}
\newunicodechar{±}{\ensuremath{\pm}}
\newunicodechar{β}{\ensuremath{\beta}}
\newunicodechar{λ}{\ensuremath{\lambda}}
\newunicodechar{σ}{\ensuremath{\sigma}}

\usepackage{xcolor}
\definecolor{rnavy}{HTML}{1B365D}
\definecolor{rgrey}{HTML}{7A828C}
\definecolor{rline}{HTML}{BFC5CC}
\definecolor{raccent}{HTML}{C0392B}

\usepackage[a4paper,top=2.6cm,bottom=2.6cm,left=2.7cm,right=2.7cm,headheight=14pt]{geometry}
\linespread{1.45}
\setlength{\parskip}{0.35em}

% 标题层级：章（\section）、节（\subsection）、小节（\subsubsection，少用）
\ctexset{
  section={format=\sffamily\Large\bfseries\color{rnavy}\raggedright, name={,}, number=\arabic{section},
           aftername=\enspace, beforeskip=1.6em, afterskip=0.8em,
           numberformat={\color{white}\colorbox{rnavy}{\makebox[1.1em]{\arabic{section}}}}},
  subsection={format=\sffamily\large\bfseries\color{rnavy}\raggedright, beforeskip=1.2em, afterskip=0.5em},
  subsubsection={format=\sffamily\normalsize\bfseries\raggedright, beforeskip=0.9em, afterskip=0.3em},
}

% 页眉：左栏目、右日期，下方细色线；页脚：页码
\usepackage{fancyhdr}
\providecommand{\ReportHeadLeft}{策略研究报告}
\providecommand{\ReportHeadRight}{2026 年 11 月}
\pagestyle{fancy}
\fancyhf{}
\fancyhead[L]{\sffamily\small\color{rnavy}\ReportHeadLeft}
\fancyhead[R]{\sffamily\small\color{rnavy}\ReportHeadRight}
\renewcommand{\headrule}{{\color{rnavy}\hrule height 0.8pt}}
\fancyfoot[C]{\sffamily\small\thepage}
\fancypagestyle{plain}{\fancyhf{}\renewcommand{\headrulewidth}{0pt}\fancyfoot[C]{\sffamily\small\thepage}}

% 图表：标题在上，「图 1　标题」
\usepackage{graphicx,float}
\usepackage[font=small,labelfont={sf,bf,color=rnavy},textfont={sf,bf,color=rnavy},labelsep=colon,justification=raggedright,singlelinecheck=false,skip=5pt]{caption}
\DeclareCaptionLabelSeparator{zhcolon}{：}
\captionsetup{labelsep=zhcolon}
\captionsetup[figure]{name=图表,position=top}
\captionsetup[table]{name=图表,position=top}
\makeatletter\let\c@table\c@figure\makeatother
% 图表下方一行资料来源（唯一允许的小字）
\newcommand{\source}[1]{\par\vspace{2pt}{\raggedright\sffamily\footnotesize\color{rgrey}资料来源：#1\par}}

\usepackage{booktabs,tabularx,array,multirow}
\renewcommand{\arraystretch}{1.3}
\newcolumntype{L}{>{\raggedright\arraybackslash}X}
\newcolumntype{C}{>{\centering\arraybackslash}X}

\usepackage{amsmath,amssymb,bm}
\usepackage{enumitem}
\setlist{nosep,leftmargin=2em,itemsep=0.2em}

\usepackage{tikz}
\usetikzlibrary{arrows.meta,positioning,calc,fit,decorations.pathreplacing}
\tikzset{
  node/.style = {draw=rnavy, line width=0.6pt, rounded corners=1.5pt, align=flush center,
                 inner xsep=6pt, inner ysep=4pt, font=\sffamily\small},
  key/.style  = {node, draw=raccent, line width=0.8pt},
  soft/.style = {node, draw=rline},
  flow/.style = {-{Stealth[length=2mm]}, line width=0.6pt, draw=rnavy},
  note/.style = {font=\sffamily\small, text=rgrey, align=flush center},
}

% 首页研究要点：每条以一句加粗结论开头
\newenvironment{keypoints}{%
  \par\noindent{\sffamily\large\bfseries\color{rnavy}研究要点}\par\vspace{0.3em}
  {\color{rnavy}\hrule height 0.6pt}\vspace{0.6em}
  \begin{itemize}[leftmargin=1.2em,itemsep=0.6em,label={\color{rnavy}\rule[0.3ex]{0.45em}{0.45em}}]}
  {\end{itemize}\vspace{0.4em}{\color{rline}\hrule height 0.4pt}}
\newcommand{\kp}[1]{{\sffamily\bfseries\color{rnavy}#1}}

\usepackage[hidelinks]{hyperref}
\usepackage{bookmark}

% 封面
\makeatletter
\newcommand{\subtitle}[1]{\gdef\@subtitle{#1}}
\newcommand{\team}[1]{\gdef\@team{#1}}
\newcommand{\reportdate}[1]{\gdef\@reportdate{#1}}
\providecommand{\@subtitle}{}\providecommand{\@team}{}\providecommand{\@reportdate}{}
\newcommand{\makereporttitle}{%
  \thispagestyle{fancy}
  \vspace*{0.4em}
  {\raggedright\sffamily\color{rnavy}\fontsize{20pt}{27pt}\selectfont\bfseries\@title\par}
  \vspace{0.5em}
  {\raggedright\sffamily\large\color{rnavy}\@subtitle\par}
  \vspace{0.3em}
  {\raggedright\sffamily\small\color{rgrey}\@team\qquad\@reportdate\par}
  \vspace{1.2em}}
\makeatother

%   \title{…} \subtitle{…} \team{…} \reportdate{…}
%   \begin{document} \makereporttitle ...

\setCJKmainfont{Songti SC}[BoldFont={Songti SC Bold}]
\def\yhpath{/Applications/Microsoft PowerPoint.app/Contents/Resources/DFonts/}
\setCJKsansfont{msyh.ttc}[Path=\yhpath,BoldFont=msyhbd.ttc]
\setCJKmonofont{msyh.ttc}[Path=\yhpath]
\setmainfont{Times New Roman}
\setsansfont{msyh.ttc}[Path=\yhpath,BoldFont=msyhbd.ttc]
\setmonofont{Menlo}[Scale=0.88]

\usepackage{newunicodechar}
\newunicodechar{→}{\ensuremath{\rightarrow}}
\newunicodechar{≈}{\ensuremath{\approx}}
\newunicodechar{≤}{\ensuremath{\leq}}
\newunicodechar{≥}{\ensuremath{\geq}}
\newunicodechar{−}{\ensuremath{-}}
\newunicodechar{×}{\ensuremath{\times}}
\newunicodechar{÷}{\ensuremath{\div}}
\newunicodechar{±}{\ensuremath{\pm}}
\newunicodechar{β}{\ensuremath{\beta}}
\newunicodechar{λ}{\ensuremath{\lambda}}
\newunicodechar{σ}{\ensuremath{\sigma}}

\usepackage{xcolor}
\definecolor{rnavy}{HTML}{1B365D}
\definecolor{rgrey}{HTML}{7A828C}
\definecolor{rline}{HTML}{BFC5CC}
\definecolor{raccent}{HTML}{C0392B}

\usepackage[a4paper,top=2.6cm,bottom=2.6cm,left=2.7cm,right=2.7cm,headheight=14pt]{geometry}
\linespread{1.45}
\setlength{\parskip}{0.35em}

% 标题层级：章（\section）、节（\subsection）、小节（\subsubsection，少用）
\ctexset{
  section={format=\sffamily\Large\bfseries\color{rnavy}\raggedright, name={,}, number=\arabic{section},
           aftername=\enspace, beforeskip=1.6em, afterskip=0.8em,
           numberformat={\color{white}\colorbox{rnavy}{\makebox[1.1em]{\arabic{section}}}}},
  subsection={format=\sffamily\large\bfseries\color{rnavy}\raggedright, beforeskip=1.2em, afterskip=0.5em},
  subsubsection={format=\sffamily\normalsize\bfseries\raggedright, beforeskip=0.9em, afterskip=0.3em},
}

% 页眉：左栏目、右日期，下方细色线；页脚：页码
\usepackage{fancyhdr}
\providecommand{\ReportHeadLeft}{策略研究报告}
\providecommand{\ReportHeadRight}{2026 年 11 月}
\pagestyle{fancy}
\fancyhf{}
\fancyhead[L]{\sffamily\small\color{rnavy}\ReportHeadLeft}
\fancyhead[R]{\sffamily\small\color{rnavy}\ReportHeadRight}
\renewcommand{\headrule}{{\color{rnavy}\hrule height 0.8pt}}
\fancyfoot[C]{\sffamily\small\thepage}
\fancypagestyle{plain}{\fancyhf{}\renewcommand{\headrulewidth}{0pt}\fancyfoot[C]{\sffamily\small\thepage}}

% 图表：标题在上，「图 1　标题」
\usepackage{graphicx,float}
\usepackage[font=small,labelfont={sf,bf,color=rnavy},textfont={sf,bf,color=rnavy},labelsep=colon,justification=raggedright,singlelinecheck=false,skip=5pt]{caption}
\DeclareCaptionLabelSeparator{zhcolon}{：}
\captionsetup{labelsep=zhcolon}
\captionsetup[figure]{name=图表,position=top}
\captionsetup[table]{name=图表,position=top}
\makeatletter\let\c@table\c@figure\makeatother
% 图表下方一行资料来源（唯一允许的小字）
\newcommand{\source}[1]{\par\vspace{2pt}{\raggedright\sffamily\footnotesize\color{rgrey}资料来源：#1\par}}

\usepackage{booktabs,tabularx,array,multirow}
\renewcommand{\arraystretch}{1.3}
\newcolumntype{L}{>{\raggedright\arraybackslash}X}
\newcolumntype{C}{>{\centering\arraybackslash}X}

\usepackage{amsmath,amssymb,bm}
\usepackage{enumitem}
\setlist{nosep,leftmargin=2em,itemsep=0.2em}

\usepackage{tikz}
\usetikzlibrary{arrows.meta,positioning,calc,fit,decorations.pathreplacing}
\tikzset{
  node/.style = {draw=rnavy, line width=0.6pt, rounded corners=1.5pt, align=flush center,
                 inner xsep=6pt, inner ysep=4pt, font=\sffamily\small},
  key/.style  = {node, draw=raccent, line width=0.8pt},
  soft/.style = {node, draw=rline},
  flow/.style = {-{Stealth[length=2mm]}, line width=0.6pt, draw=rnavy},
  note/.style = {font=\sffamily\small, text=rgrey, align=flush center},
}

% 首页研究要点：每条以一句加粗结论开头
\newenvironment{keypoints}{%
  \par\noindent{\sffamily\large\bfseries\color{rnavy}研究要点}\par\vspace{0.3em}
  {\color{rnavy}\hrule height 0.6pt}\vspace{0.6em}
  \begin{itemize}[leftmargin=1.2em,itemsep=0.6em,label={\color{rnavy}\rule[0.3ex]{0.45em}{0.45em}}]}
  {\end{itemize}\vspace{0.4em}{\color{rline}\hrule height 0.4pt}}
\newcommand{\kp}[1]{{\sffamily\bfseries\color{rnavy}#1}}

\usepackage[hidelinks]{hyperref}
\usepackage{bookmark}

% 封面
\makeatletter
\newcommand{\subtitle}[1]{\gdef\@subtitle{#1}}
\newcommand{\team}[1]{\gdef\@team{#1}}
\newcommand{\reportdate}[1]{\gdef\@reportdate{#1}}
\providecommand{\@subtitle}{}\providecommand{\@team}{}\providecommand{\@reportdate}{}
\newcommand{\makereporttitle}{%
  \thispagestyle{fancy}
  \vspace*{0.4em}
  {\raggedright\sffamily\color{rnavy}\fontsize{20pt}{27pt}\selectfont\bfseries\@title\par}
  \vspace{0.5em}
  {\raggedright\sffamily\large\color{rnavy}\@subtitle\par}
  \vspace{0.3em}
  {\raggedright\sffamily\small\color{rgrey}\@team\qquad\@reportdate\par}
  \vspace{1.2em}}
\makeatother

%   \title{…} \subtitle{…} \team{…} \reportdate{…}
%   \begin{document} \makereporttitle ...

\setCJKmainfont{Songti SC}[BoldFont={Songti SC Bold}]
\def\yhpath{/Applications/Microsoft PowerPoint.app/Contents/Resources/DFonts/}
\setCJKsansfont{msyh.ttc}[Path=\yhpath,BoldFont=msyhbd.ttc]
\setCJKmonofont{msyh.ttc}[Path=\yhpath]
\setmainfont{Times New Roman}
\setsansfont{msyh.ttc}[Path=\yhpath,BoldFont=msyhbd.ttc]
\setmonofont{Menlo}[Scale=0.88]

\usepackage{newunicodechar}
\newunicodechar{→}{\ensuremath{\rightarrow}}
\newunicodechar{≈}{\ensuremath{\approx}}
\newunicodechar{≤}{\ensuremath{\leq}}
\newunicodechar{≥}{\ensuremath{\geq}}
\newunicodechar{−}{\ensuremath{-}}
\newunicodechar{×}{\ensuremath{\times}}
\newunicodechar{÷}{\ensuremath{\div}}
\newunicodechar{±}{\ensuremath{\pm}}
\newunicodechar{β}{\ensuremath{\beta}}
\newunicodechar{λ}{\ensuremath{\lambda}}
\newunicodechar{σ}{\ensuremath{\sigma}}

\usepackage{xcolor}
\definecolor{rnavy}{HTML}{1B365D}
\definecolor{rgrey}{HTML}{7A828C}
\definecolor{rline}{HTML}{BFC5CC}
\definecolor{raccent}{HTML}{C0392B}

\usepackage[a4paper,top=2.6cm,bottom=2.6cm,left=2.7cm,right=2.7cm,headheight=14pt]{geometry}
\linespread{1.45}
\setlength{\parskip}{0.35em}

% 标题层级：章（\section）、节（\subsection）、小节（\subsubsection，少用）
\ctexset{
  section={format=\sffamily\Large\bfseries\color{rnavy}\raggedright, name={,}, number=\arabic{section},
           aftername=\enspace, beforeskip=1.6em, afterskip=0.8em,
           numberformat={\color{white}\colorbox{rnavy}{\makebox[1.1em]{\arabic{section}}}}},
  subsection={format=\sffamily\large\bfseries\color{rnavy}\raggedright, beforeskip=1.2em, afterskip=0.5em},
  subsubsection={format=\sffamily\normalsize\bfseries\raggedright, beforeskip=0.9em, afterskip=0.3em},
}

% 页眉：左栏目、右日期，下方细色线；页脚：页码
\usepackage{fancyhdr}
\providecommand{\ReportHeadLeft}{策略研究报告}
\providecommand{\ReportHeadRight}{2026 年 11 月}
\pagestyle{fancy}
\fancyhf{}
\fancyhead[L]{\sffamily\small\color{rnavy}\ReportHeadLeft}
\fancyhead[R]{\sffamily\small\color{rnavy}\ReportHeadRight}
\renewcommand{\headrule}{{\color{rnavy}\hrule height 0.8pt}}
\fancyfoot[C]{\sffamily\small\thepage}
\fancypagestyle{plain}{\fancyhf{}\renewcommand{\headrulewidth}{0pt}\fancyfoot[C]{\sffamily\small\thepage}}

% 图表：标题在上，「图 1　标题」
\usepackage{graphicx,float}
\usepackage[font=small,labelfont={sf,bf,color=rnavy},textfont={sf,bf,color=rnavy},labelsep=colon,justification=raggedright,singlelinecheck=false,skip=5pt]{caption}
\DeclareCaptionLabelSeparator{zhcolon}{：}
\captionsetup{labelsep=zhcolon}
\captionsetup[figure]{name=图表,position=top}
\captionsetup[table]{name=图表,position=top}
\makeatletter\let\c@table\c@figure\makeatother
% 图表下方一行资料来源（唯一允许的小字）
\newcommand{\source}[1]{\par\vspace{2pt}{\raggedright\sffamily\footnotesize\color{rgrey}资料来源：#1\par}}

\usepackage{booktabs,tabularx,array,multirow}
\renewcommand{\arraystretch}{1.3}
\newcolumntype{L}{>{\raggedright\arraybackslash}X}
\newcolumntype{C}{>{\centering\arraybackslash}X}

\usepackage{amsmath,amssymb,bm}
\usepackage{enumitem}
\setlist{nosep,leftmargin=2em,itemsep=0.2em}

\usepackage{tikz}
\usetikzlibrary{arrows.meta,positioning,calc,fit,decorations.pathreplacing}
\tikzset{
  node/.style = {draw=rnavy, line width=0.6pt, rounded corners=1.5pt, align=flush center,
                 inner xsep=6pt, inner ysep=4pt, font=\sffamily\small},
  key/.style  = {node, draw=raccent, line width=0.8pt},
  soft/.style = {node, draw=rline},
  flow/.style = {-{Stealth[length=2mm]}, line width=0.6pt, draw=rnavy},
  note/.style = {font=\sffamily\small, text=rgrey, align=flush center},
}

% 首页研究要点：每条以一句加粗结论开头
\newenvironment{keypoints}{%
  \par\noindent{\sffamily\large\bfseries\color{rnavy}研究要点}\par\vspace{0.3em}
  {\color{rnavy}\hrule height 0.6pt}\vspace{0.6em}
  \begin{itemize}[leftmargin=1.2em,itemsep=0.6em,label={\color{rnavy}\rule[0.3ex]{0.45em}{0.45em}}]}
  {\end{itemize}\vspace{0.4em}{\color{rline}\hrule height 0.4pt}}
\newcommand{\kp}[1]{{\sffamily\bfseries\color{rnavy}#1}}

\usepackage[hidelinks]{hyperref}
\usepackage{bookmark}

% 封面
\makeatletter
\newcommand{\subtitle}[1]{\gdef\@subtitle{#1}}
\newcommand{\team}[1]{\gdef\@team{#1}}
\newcommand{\reportdate}[1]{\gdef\@reportdate{#1}}
\providecommand{\@subtitle}{}\providecommand{\@team}{}\providecommand{\@reportdate}{}
\newcommand{\makereporttitle}{%
  \thispagestyle{fancy}
  \vspace*{0.4em}
  {\raggedright\sffamily\color{rnavy}\fontsize{20pt}{27pt}\selectfont\bfseries\@title\par}
  \vspace{0.5em}
  {\raggedright\sffamily\large\color{rnavy}\@subtitle\par}
  \vspace{0.3em}
  {\raggedright\sffamily\small\color{rgrey}\@team\qquad\@reportdate\par}
  \vspace{1.2em}}
\makeatother

\title{端到端赛道 AI 工具使用说明}
\subtitle{分钟级 A 股截面预测 · 端到端赛道决赛提交材料}
\team{fx\hbox{二号机}}
\reportdate{2026 年 11 月}

% 本文件内的排版设置：表格行距收紧，减少只放浮动体的页面
\newcommand{\tabsetup}{\small\linespread{1.15}\selectfont\renewcommand{\arraystretch}{1.15}}
\renewcommand{\topfraction}{0.9}
\renewcommand{\bottomfraction}{0.8}
\renewcommand{\textfraction}{0.07}
\renewcommand{\floatpagefraction}{0.85}
\setcounter{topnumber}{3}
\setcounter{totalnumber}{4}

\begin{document}
\makereporttitle

本作品使用命令行编程代理 Claude Code（下文称 AI 编码代理）编写代码、设计并运行实验、分析结果、撰写记录，并起草决赛材料；评价标准、合规判断、平台操作、提交决策和私榜作品选择由团队成员决定。团队对全部提交内容负责。

\section{使用的工具与版本}

使用的 AI 工具是 Anthropic 公司的 Claude Code，模型为 claude-opus-5-5，客户端版本为 2.1.293～2.1.321。2026 年 10 月 10 日至 28 日用于开发与提交，10 月 31 日用于收尾，11 月用于整理决赛材料。该工具在团队成员的本地计算机上运行，经授权通过远程登录在两台训练机的指定目录内读写文件、运行训练，不接触比赛平台。端到端目录的 205 个代码提交中有 204 个由其提交，唯一例外是仓库初始化提交，其中收录的项目启动说明与早期交接文档写于 10 月 10 日之前，作者身份未留存记录。因子赛道使用过的批量编码代理没有参与端到端赛道。

\section{使用环节}

图表 1 按研究环节列出 AI 编码代理与团队成员各自承担的工作。

\begin{table}[!htbp]
\caption{各研究环节中 AI 编码代理与团队成员承担的工作}
\centering\tabsetup
\begin{tabularx}{\textwidth}{@{}p{5em}>{\raggedright\arraybackslash\hsize=1.1\hsize}X>{\raggedright\arraybackslash\hsize=0.9\hsize}X@{}}
\toprule
研究环节 & AI 编码代理承担的工作 & 团队成员承担的工作\\
\midrule
评测口径 & 编写平台探测脚本，核实官方标签定义；本地评测改为平台口径 & 在平台运行探测脚本并带回结果；转达主办方规则与答复\\
输入设计 & 编写数据读取、切窗、归一化与特征计算代码，加入特征时点断言 & 逐项追问新输入是否合规，否决合规存疑的方案\\
模型设计 & 设计三尺度网络、同伴注意力（基于走势相似股票的注意力机制）、秩高斯目标、多空组合换手正则项，以及梯度提升树与 GRU 子模型 & 听取设计说明后才开始训练；提出研究方向，如关注子模型间的信息差异\\
实验与分析 & 在两台训练机上排队运行留出验证实验，每次只改一个变量；分年份按平台口径比较五项指标 & 规定训练串行、日志可见；确定全量训练准入标准；追问收益来源\\
公榜验证 & 设计单变量线上验证，事先写下判据；记录公榜结果 & 决定是否占用当日额度；在平台上传并回传结果\\
打包与交付 & 编写打包、校验和截断检测脚本；已交付版本不再修改 & 决定提交与否，选择私榜作品\\
决赛材料 & 依据仓库记录起草研究报告、展示材料、讲稿与本说明 & 审阅全部材料，提出修改意见，确定最终文本\\
\bottomrule
\end{tabularx}
\source{团队整理}
\end{table}

两个赛道共用同一个 AI 编码代理，模型输出与因子值互相隔离：端到端模型的输入里没有因子值或因子输出，因子代码里没有端到端预测；研究思路相互借鉴，例如梯度提升树子模型的部分截面特征受因子赛道研究方向启发。

\section{关键提示词示例}

决定方向、标准和取舍的提示词，主要是对话中的追问与要求。图表 2 按环节概括其中 6 条提示词的内容，均为转述。

\begin{table}[!htbp]
\caption{关键提示词的内容（按环节转述）}
\centering\tabsetup
\begin{tabularx}{\textwidth}{@{}p{4em}>{\raggedright\arraybackslash}X>{\raggedright\arraybackslash}X@{}}
\toprule
环节 & 提示词要求 & 引出的结果\\
\midrule
评测口径 & 查明本地评测与公榜差距大的原因 & 编写平台探测脚本，查出标签口径有误\\
\midrule
合规 & 检查信号是否过强、是否合规 & 编写截断检测脚本，每版打包时运行\\
\midrule
研究方法 & 五项指标都要更好，不以排名为目标 & 改动按两年分别比较五项指标\\
 & 分析收益来源，不只比较随机种子 & 转做收益来源分析，引出权重滑动平均\\
 & 先确认改动有提升，再做全量训练 & 叫停未经验证的全量训练\\
\midrule
作品选择 & 检查是否对公榜过拟合 & 含 4 个 GRU 子模型的版本不提交\\
\bottomrule
\end{tabularx}
\source{团队整理（依据会话记录转述）}
\end{table}

较长的全天工作目标由 AI 编码代理起草，经团队成员审定后设定；目标规定对话中的最新指示优先于目标本身，并要求 AI 编码代理不代为上传，也不催促上传。

\section{人工复核方式}

\subsection{复核层次}

AI 编码代理的产出在采用之前依次经过图表 3 所列的检查。其中评测口径对齐、全量训练准入标准和截断检测三项，均在团队成员质疑或追问之后建立。

\begin{table}[!htbp]
\caption{人工复核的层次与查出过的问题}
\centering\tabsetup
\begin{tabularx}{\textwidth}{@{}p{4.4em}LL@{}}
\toprule
检查 & 内容与执行方式 & 查出过的问题\\
\midrule
评测口径\newline 对齐 & 本地评测按平台口径计算五项指标，标签定义经平台探测核实 & 早期标签与官方标签秩相关只有 0.56\\
留出验证 & 2020–2022 年训练，2023、2024 年分别评估；差值超过随机种子噪声才算有效 & 多项结构改动换随机种子复现后持平，未采用\\
全量训练\newline 准入标准 & RankIC ≥~+0.005、ICIR ≥~+0.08、压力稳定性 ≥~+0.05、多空组合夏普比率 ≥~+1.5、换手率 ≤~+0.01；规则由团队成员设定，阈值由 AI 编码代理提出 & 多个已打包并通过校验的版本因增量不足未提交\\
单变量\newline 线上验证 & 只改一个变量的单模型提交公榜，与对照同一时点比较，判据事先写下 & 全天 5 分钟片段输入本地为正、公榜 RankIC −0.0033，被否决\\
交付前\newline 检查 & 打包校验；截断检测要求只用截断时刻之前数据的推理结果与使用全量数据时逐行一致 & 最终模型在两个截断时刻分别比较 19,000 行和 24,000 行，差异均为 0\\
\bottomrule
\end{tabularx}
\source{团队整理}
\end{table}

\subsection{典型纠偏}

第一个例子是评测标签。早期本地评测使用按推测构造的标签，AI 编码代理起初认为该标签正确。团队成员认为本地评分与公榜差距过大，提出在平台上代为运行探测脚本；核对后确认原标签与官方标签的截面秩相关只有 0.56，修正后约为 0.97，此前的本地结果全部作废。

第二个例子是合规判断。一次诊断显示传统因子在公榜上明显好于当时的神经网络，AI 编码代理据此建议以传统因子为基础。团队成员追问规则是否允许这样做，并决定采取稳妥路线：端到端模型只使用规则允许的输入，输出由训练学习得到。

第三个例子是实验流程。多空组合换手正则项强度 0.3 的单变量线上验证结果尚未返回，AI 编码代理已安排该强度的全量训练；结果返回后 ICIR 由强度 0.1 时的 1.990 降至 1.786，它又将强度改为 0.15，仍未经验证。经团队成员追问，这次全量训练被取消。

复核也有遗漏。对照冻结文件发现，预留的四个 30 分钟窗口盘口均值通道在冻结权重中恒为 0，不参与预测；打包校验和截断检测均未覆盖这一点，AI 编码代理此前对模型输入的描述也与之不符。推理输出不受影响，研究报告按实际起作用的输入描述模型。AI 编码代理对自身代码的描述同样需要对照冻结文件复核。

\section{最终采用情况}

比赛期间，AI 编码代理完成至少 223 次留出验证训练和 39 次全量训练，打包 24 个正式提交版本；有公榜结果的提交 38 个，进入私榜的只有 1 个。

最终模型由 4 个 Transformer、2 个 GRU 和 1 个梯度提升树等权集成，各组成部分均由 AI 编码代理设计实现。09:30 截面输出置零和训练目标对行情代理风格因子与聚类行业取残差，起因是团队成员转达的规则和平台探测结果，这两项连同推理端风格中性化合计使公榜压力稳定性提高 0.16。当前时段输入源于团队成员对收益来源的追问，单变量线上验证中 ICIR 提高 0.089；秩高斯目标在单变量线上验证中使 ICIR 提高 0.193；多空组合换手正则项强度为 0.1 时，换手率降低 0.036。2 个 GRU 子模型经团队成员批准后试验，公榜换手率降低 0.006，其余四项持平。

被否决的 AI 产出同样有记录。除第 4.2 节的例子外，推理时对模型自身之前的输出做平滑违反规则，当天撤回；加宽、加深等结构尝试未通过本地验证；含 4 个 GRU 子模型的版本在四个半年的 ICIR 都略低，没有提交。

私榜每队只能选一个作品。团队成员选择了相对上一版集成只降低换手率、其余四项持平的版本，这一点在本地两年、四个半年和公榜同一时点比较中一致。该作品在私榜排名第六，私榜成绩由平台从零重训的模型评出。

\section{责任声明}

AI 编码代理不是本作品的作者，其产出经第 4 节所列方式复核后才被采用。作品中的代码、实验结果与全部材料由团队负责，数字与结论经团队成员审阅。

AI 编码代理不直接登录或提交比赛平台：作品上传、探测脚本运行、公榜结果回传和私榜作品选择均由团队成员完成，删除文件等不可逆操作也由团队成员执行。代码提交历史、实验日志、提交记录和完整对话记录均已留存，可应主办方要求提供。

\end{document}
